% SPDX-FileCopyrightText: 2026 Will Cook
% SPDX-License-Identifier: CC-BY-4.0
% Open-source strategy paper; build with tectonic.
\documentclass[10pt]{article}
\usepackage{paper-house-style}
\usepackage{booktabs,array,float,placeins,multicol,graphicx}
\usepackage{tikz}
\usetikzlibrary{arrows.meta,positioning,calc,fit,shapes.geometric,matrix}
\usepackage[colorlinks=true,linkcolor=linkinternal,urlcolor=linkexternal,citecolor=linkinternal]{hyperref}
\hypersetup{
  bookmarksnumbered=true,bookmarksopen=true,bookmarksopenlevel=1,
  pdftitle={From Spare Compute to Cumulative Mathematics},
  pdfauthor={Will Cook},
  pdfsubject={An open-source strategy for agent-assisted mathematical research},
  pdflang={en-GB},
  pdfkeywords={open source, formal mathematics, Lean 4, AI agents, volunteer computing, research commons, attribution},
}
\newcommand{\repobase}{https://github.com/wcook04/plectis-erdos}
\newcommand{\repolink}[2]{\href{\repobase/blob/ca0e13f8acf5ccf48506e4bdb870953d3a0856fa/#1}{\texttt{#2}}}
\newcommand{\repoissue}[2]{\href{\repobase/issues/new?template=#1}{#2}}
\emergencystretch=1.7em

\runninghead{From spare compute to cumulative mathematics}
\title{From Spare Compute to Cumulative Mathematics}
\subtitle{An open-source strategy for agent-assisted research on hard problems}
\author{Will Cook}
\date{September 2026}

\begin{document}
\maketitle
\paperprovenance
\fontsize{9.2pt}{11pt}\selectfont

\begin{abstract}
We want someone with a mathematical idea, spare compute, or time to review a
proof to be able to contribute to an ongoing research problem without first
building a research environment. Plectis puts the papers, formal source,
experiments, and tools in a public Git repository. A contributor can clone it,
work on one question with an agent or by hand, and return a result for review.
Accepted contributions retain their evidence and credit.

The repository currently covers eight open Erd\H{o}s problems. A useful return
may be a proof, but it may also be a counterexample, a corrected reference, an
explanation of why a method fails, or a repair to the tools. The proposal is
that these contributions should accumulate together: later attempts should
have both better mathematics to draw on and a better environment for doing
the work. This paper describes the contribution process, its precedents, and
the tests needed to assess it. The public workflow is implemented, but its
usefulness to outside contributors and its
effect on research productivity have not yet been established.
\end{abstract}

\section{The strategy}\label{sec:strategy}

The project began with one undergraduate, without an institutional research
team or a dedicated compute allocation. Preparing an agent to do useful work
required more than choosing a model. It required locating the relevant
mathematics, setting up proof checks, keeping track of unsuccessful attempts,
and writing an account that could be understood after the conversation ended.
The reason for publishing the environment is to let others use and improve
that work.

A contributor need not take on a whole problem. A mathematician may suggest a
missing lemma or find a counterexample. A programmer may run an exact
calculation. A Lean contributor may formalise an argument or discover that its
formal statement is wrong. Reviewers and expositors may supply what a long
agent run could not: a sound judgement of the result and an explanation of its
main idea. Contributions to the software are welcome on the same terms.

The eight Erd\H{o}s problems give this work a mathematical purpose. A run can
leave one of them open and still establish something worth preserving.
Tao's discussion of learning from the study of a problem is relevant here
\cite{taolearning2026}: a failed construction may reveal an invariant, expose
a limitation of a technique, or suggest a connection elsewhere. Such a result
needs an account of why the approach seemed plausible and where it stopped.
His Navier--Stokes example raises the converse concern: publishing a final
construction while concealing its development may lose much of the insight
that led to it \cite{taoexploration2026}.

We therefore publish selected research records while the problems are still
open. The records include proved statements, computations, references,
counterexamples, and unresolved steps. An agent retrieves the relevant work,
saves its attempted proofs and check results, and updates the affected files
when a result is accepted. The
\papersectionlink{claim-faithful-publication-systems-paper.pdf}%
{systems-research-attempt}{systems paper's composite-dilation example}
follows an unsuccessful extension through this process.

\subsection*{Research and review}
\papersectiontarget{strategy-coupled-goals}

The public workflow has two recurring jobs. The first investigates a
mathematical question. The second checks how the result fits the existing
work and updates the paper and repository. The skill files call these
\emph{discovery} and \emph{stewardship}. One agent can alternate between them;
contributors can also divide the work.

The second job is needed because a successful proof attempt does not settle
how the result should be presented or used. Five new declarations may amount
to one routine lemma. A short counterexample may eliminate a line of research
that was consuming most of the compute. A literature search may uncover a
stronger theorem than the one leading the paper. The review changes the
account of the result and, when appropriate, the next question to investigate.

The \repolink{skills/run-coupled-research-goals/SKILL.md}%
{coupled-research skill} passes source revisions, results, and check records
between these jobs. A new theorem, correction, or review can trigger another
pass. When nothing has changed, the jobs yield. The
\papersectionlink{claim-faithful-publication-systems-paper.pdf}%
{systems-coupled-goals}{systems paper} explains the implementation; the
\papersectionlink{cold-clone-to-proof-receipt.pdf}%
{cold-clone-problem}{navigation paper} examines how a new agent finds the
mathematics and its checks. All of these public jobs can run without the
private workbench.

\subsection*{What one clone lets a contributor do}

A newcomer can ask an agent to explain the repository before choosing a task.
The \repolink{skills/explain-public-system/SKILL.md}%
{explanation skill} directs it to the public papers and source, with the level
of detail adapted to the reader. Figure~\ref{fig:first-actions} gives examples
of work that can follow. A contributor can also read a problem paper directly
and proceed without an agent.

\begin{figure}[H]
\centering
\begin{tikzpicture}[
  action/.style={paper stage,text width=3.15cm,minimum height=15mm},
  arr/.style={paper arrow}
]
\node[action,fill=MidnightBlue!5] (mine)
  {{\scriptsize\sffamily\bfseries 1}\\[2pt]\textbf{Investigate a question}\\[2pt]
   \scriptsize bounded theorem or obstruction};
\node[action,fill=MidnightBlue!5,right=6mm of mine] (direct)
  {{\scriptsize\sffamily\bfseries 2}\\[2pt]\textbf{Give direction}\\[2pt]
   \scriptsize hypothesis, example, or counterexample};
\node[action,fill=TealBlue!5,right=6mm of direct] (formal)
  {{\scriptsize\sffamily\bfseries 3}\\[2pt]\textbf{Check a claim}\\[2pt]
   \scriptsize formalise or review};
\node[action,fill=BurntOrange!6,below=6mm of $(mine.south)!0.5!(direct.south)$] (repair)
  {{\scriptsize\sffamily\bfseries 4}\\[2pt]\textbf{Repair the machinery}\\[2pt]
   \scriptsize navigation, validation, or exposition};
\node[action,fill=BurntOrange!6,below=6mm of $(direct.south)!0.5!(formal.south)$] (add)
  {{\scriptsize\sffamily\bfseries 5}\\[2pt]\textbf{Add a problem}\\[2pt]
   \scriptsize sourced world and boundary};
\node[paper box,below=8mm of $(repair.south)!0.5!(add.south)$,
      text width=10.8cm,fill=ForestGreen!6,align=center] (return)
      {{\scriptsize\sffamily\bfseries COMMON OUTPUT}\\[3pt]
       \textbf{Result with evidence}\\[2pt]
       \footnotesize source \(\cdot\) evidence \(\cdot\) limitation \(\cdot\) open statement};
\draw[arr] (mine.south) -- (repair.north);
\draw[arr] (direct.south) -- ($(repair.north)!0.5!(add.north)$);
\draw[arr] (formal.south) -- (add.north);
\draw[arr] (repair.south) -- (repair.south |- return.north);
\draw[arr] (add.south) -- (add.south |- return.north);
\end{tikzpicture}
\caption{Ways to contribute, from investigating a question to improving the tools.}
\label{fig:first-actions}
\end{figure}

\section{Where further work is needed}\label{sec:model}

Research throughput depends on more than the number of agents running.
Mathematical direction determines which questions they attempt; the available
literature and tools determine what they can build on; verification and review
determine which outputs become usable results. These constraints suggest two
ways to contribute. A theorem, counterexample, or correction changes the
mathematical record. Better retrieval, validation, or exposition changes the
cost of using that record in later work.

For example, an unsuccessful extension may reveal a missing hypothesis.
Recording that obstruction changes the next mathematical attempt. If the
attempt also exposed a stale source link, repairing the link improves the
environment for every subsequent reader. The two changes have different
checks and can be adopted and credited separately. Figure~\ref{fig:conversion}
shows how contributions enter this cycle. Its effect on productivity is an
empirical question for Section~\ref{sec:evaluation}.

\begin{figure}[H]
\centering
\begin{tikzpicture}[
  column/.style={paper box,align=center,text width=3.15cm,minimum height=5.0cm},
  receipt/.style={paper box,align=left,text width=5.05cm,minimum height=1.5cm,
                  font=\footnotesize},
  arr/.style={paper arrow}
]
\node[column,fill=MidnightBlue!5] (roles)
  {{\scriptsize\sffamily\bfseries OUTSIDE FUNCTIONS}\\[4pt]
   {\large\textbf{Inputs}}\\[6pt]
   \footnotesize mathematical direction\\[2pt]
   compute and agents\\[2pt]
   formalisation\\[2pt]
   infrastructure\\[2pt]
   review and exposition\\[2pt]
   new problems and funding};
\node[column,fill=TealBlue!5,right=8mm of roles] (engine)
  {{\scriptsize\sffamily\bfseries SHARED ENGINE}\\[4pt]
   {\large\textbf{Candidate production}}\\[6pt]
   \footnotesize \textbf{1} choose a frontier\\[3pt]
   \textbf{2} read the relevant corpus\\[3pt]
   \textbf{3} reason, compute, formalise\\[3pt]
   \textbf{4} return a candidate};
\node[column,fill=ForestGreen!6,right=8mm of engine] (result)
  {{\scriptsize\sffamily\bfseries PUBLIC RESULT}\\[4pt]
   {\large\textbf{Review funnel}}\\[6pt]
   \footnotesize replay and challenge\\[3pt]
   statement reconciliation\\[3pt]
   \rule{2.3cm}{.35pt}\\[5pt]
   \textbf{Accepted corpus}\\[2pt]
   public provenance};
\draw[arr] (roles) -- node[paper label,above]{contributions} (engine);
\draw[arr] (engine) -- node[paper label,above]{candidate} (result);

\node[receipt,below=7mm of $(roles.south)!0.5!(engine.south)$] (resource)
  {{\scriptsize\sffamily\bfseries RESOURCE RECEIPT}\\[2pt]
   human time; compute; model and harness; starting commit};
\node[receipt,right=7mm of resource] (lineage)
  {{\scriptsize\sffamily\bfseries CONTRIBUTION RECEIPT}\\[2pt]
   who supplied what; supporting evidence; later uses and repairs};
\node[paper box,below=6mm of $(resource.south)!0.5!(lineage.south)$,
      text width=10.9cm,align=center,fill=BurntOrange!6] (feedback)
  {Accepted mathematics enlarges the corpus; accepted architecture improves later runs.};
\draw[paper feedback] (feedback.west) to[out=180,in=-90] (roles.south);
\end{tikzpicture}
\caption{Contributors supply ideas, computation, formalisation, software, and review.
The accepted work updates the mathematics or the tools used in later attempts.}
\label{fig:conversion}
\end{figure}

A versioned repository also permits longitudinal evaluation.  A fixed,
versioned problem world can be revisited as models, runners, and compute
budgets change.  Comparisons must disclose the starting commit, model,
scaffolding, number of attempts, compute, and review procedure.  Otherwise a
success says little about which variable changed.  A public failure record is
equally important: capability claims are badly distorted when successes are
announced and the number and cost of failed attempts are hidden
\cite{tao2026}.

\section{Why begin with hard mathematics}\label{sec:math}

Mathematical work can travel with the tools needed to examine it. A clone can
contain papers, exact computations, and formal proofs, and a contributor can
check one lemma without first completing the entire argument. This makes
mathematics a practical starting point for the proposal.

Computation helps choose what to try. It can find a counterexample, compare
two formulations, or reveal a pattern that suggests a proof. The record needs
the inputs, code, output, and interpretation, so a later reader can distinguish
what was calculated from what was conjectured.

The eight selected problems provide different kinds of work
(Figure~\ref{fig:problem-worlds}). Their value as research environments will
depend on the questions and results they produce, rather than on their
recognisable problem numbers.

\begin{figure}[H]
\centering
\begin{tikzpicture}[
  card/.style={paper card,text width=5.05cm,minimum height=1.35cm}
]
\matrix[column sep=5mm,row sep=3mm] {
  \node[card,fill=MidnightBlue!4] {\textbf{68 \quad Factorial reciprocal series}\\[-1pt]
    \footnotesize carries and denominator escape}; &
  \node[card,fill=MidnightBlue!4] {\textbf{243 \quad Reciprocal-tail rigidity}\\[-1pt]
    \footnotesize discrete dynamics and recovery}; \\
  \node[card,fill=TealBlue!4] {\textbf{249 \quad Binary totient series}\\[-1pt]
    \footnotesize arithmetic tails and anti-concentration}; &
  \node[card,fill=TealBlue!4] {\textbf{251 \quad Prime-gap dyadic series}\\[-1pt]
    \footnotesize prime structure versus tail escape}; \\
  \node[card,fill=ForestGreen!5] {\textbf{257 \quad Mersenne-support subseries}\\[-1pt]
    \footnotesize achievement sets and rational targets}; &
  \node[card,fill=ForestGreen!5] {\textbf{269 \quad Running-lcm sums}\\[-1pt]
    \footnotesize finite-state structure and carries}; \\
  \node[card,fill=BurntOrange!5] {\textbf{1041 \quad Lemniscate geometry}\\[-1pt]
    \footnotesize local flow versus global path geometry}; &
  \node[card,fill=BurntOrange!5] {\textbf{1049 \quad Rational-base series}\\[-1pt]
    \footnotesize Lambert approximation kernels and obstructions}; \\
};
\end{tikzpicture}
\caption{Eight distinct research environments.  The results map and the
problem papers state their exact frontiers; this strategy paper does not
reproduce the specialist vocabulary needed to attack them.}
\label{fig:problem-worlds}
\end{figure}

Every row admits more than theorem proving.  A mathematician may sharpen a
hypothesis or supply a counterexample.  A programmer may run a discriminating
exact calculation.  A formaliser may turn a paper deduction into a checked
declaration or find that the two statements disagree.  A literature reader
may locate a theorem that closes or invalidates a route.  The contribution is
the exact returned object and its limitation, not the prestige of the problem
number.

\section{The public research object}\label{sec:object}

The unit of participation is a \emph{problem world}.  A problem world includes
the endpoint question, current public status, principal theorems, formal
source, experiments, known counterexamples, failed mechanisms, source
literature, and exact open obligations.  A bounded task is a neighbourhood in
that world, not an invitation to read the repository indiscriminately.

The public layers carry different kinds of evidence.  Their order is easier to
read as an authority ladder than as one all-purpose badge.

\begin{figure}[H]
\centering
\begin{tikzpicture}[
  rung/.style={paper stage,text width=3.2cm,minimum height=2.35cm},
  wide/.style={rung,text width=4.95cm},
  arr/.style={paper arrow}
]
\node[rung,fill=MidnightBlue!5] (exp)
  {{\scriptsize\sffamily\bfseries 1 \; FINITE EVIDENCE}\\[3pt]
   \textbf{Experiment}\\[3pt]\footnotesize establishes a finite output\\[2pt]
   \textit{Scope: no unbounded theorem}};
\node[rung,fill=TealBlue!5,right=6mm of exp] (lean)
  {{\scriptsize\sffamily\bfseries 2 \; FORMAL EVIDENCE}\\[3pt]
   \textbf{Lean kernel}\\[3pt]\footnotesize accepts the exact proposition\\[2pt]
   \textit{Scope: intended meaning is separate}};
\node[rung,fill=ForestGreen!5,right=6mm of lean] (match)
  {{\scriptsize\sffamily\bfseries 3 \; ALIGNMENT}\\[3pt]
   \textbf{Statement review}\\[3pt]\footnotesize compares formal and informal claims\\[2pt]
   \textit{Scope: no novelty judgement}};
\node[wide,fill=BurntOrange!6,below=8mm of match] (expert)
  {{\scriptsize\sffamily\bfseries 4 \; SCHOLARLY REVIEW}\\[3pt]
   \textbf{Expert assessment}\\[3pt]\footnotesize soundness, relevance, and relation to prior work};
\node[wide,fill=Violet!5,left=7mm of expert] (community)
  {{\scriptsize\sffamily\bfseries 5 \; EXTERNAL OUTCOME}\\[3pt]
   \textbf{Community uptake}\\[3pt]\footnotesize develops through criticism, use, and reuse};
\draw[arr] (exp) -- (lean);
\draw[arr] (lean) -- (match);
\draw[arr] (match) -- (expert);
\draw[arr] (expert) -- (community);
\end{tikzpicture}
\caption{Evidence and acceptance remain separate.  A candidate need not pass
through every box in a single line, but no earlier box inherits the authority
of a later one.}
\label{fig:authority-ladder}
\end{figure}

The repository uses established tools and contribution practices.  Erd\H{o}s Problems supplies questions, sources, status,
and a problem community \cite{erdosproblems}.  Lean supplies the formal
language and kernel \cite{lean4}.  DeepMind's
\texttt{formal-conjectures} supplies a public one-problem-file model with
metadata, snapshots, issues, and pull requests \cite{formalconjectures}.
Comparator and Palomar supply exact-interface checking, permitted-axiom
inspection, and a durable formal record with a deliberately limited editorial
claim \cite{palomar}.  Polymath supplies a social precedent for publishing
small, tentative, and negative contributions \cite{polymath}.  BOINC and GIMPS
supply the volunteer-compute precedent and make visible why executable work,
independent checks, and legible credit matter \cite{boinc,gimps}.  The present
repository wraps these practices into a route that a newcomer can enter
without first rebuilding each component.

The same underlying substrate can support several reader projections: a short
public primer, a specialist paper, a detailed proof account, Lean declarations,
an agent explanation, and machine-readable queries.  Each projection must link
back to its source claims and evidence.  A shorter or friendlier account does
not acquire permission to strengthen them.

An external runner begins at \repolink{AGENTS.override.md}{the compact agent
entry}.  It can inspect all problem frontiers, choose a bounded question, read
the relevant paper and source neighbourhood, run experiments or edit formal
code, validate the result, and prepare a return.  None of these steps requires
access to the private workbench.  A contributor is free to replace the agent,
the scheduler, or the entire search policy while retaining the public evidence
boundary.

\section{The contribution protocol}\label{sec:protocol}
\papersectiontarget{strategy-protocol}

A contributor forks the repository, works in a local clone, and opens a pull
request with the proposed change. A finding without a patch can be submitted
as a research-progress issue; an infrastructure proposal has its own issue
form. The \texttt{submit-pull-request} skill helps an agent prepare the files,
checks, and description. Publishing the branch and request requires the
contributor's authorisation.

A return records the starting commit, the result, the evidence, the tools used,
and the people who contributed. Proof changes need the appropriate Lean
checks. Computations need enough code and input data to reproduce the output.
A proposed strengthening must say which hypothesis or conclusion changed.
Reviewers then check the result, its relation to the stated problem and prior
work, and the changes needed in the paper.

The repository may have moved on while the contribution was being prepared.
The starting commit lets a reviewer reconstruct the original change before
adapting it to current main. Substantive conflict resolution is recorded as
integration work, preserving the original contributor's credit. The final
checks run against the version that will be adopted.

Before submission, the contributor also checks what the result affects. A new
lemma may change a later proof, a claimed open question, a computation, or a
passage in the paper. The \texttt{propagate-research-consequences} skill asks
for each affected file to be updated, checked as unchanged, or deferred with
a reason. After review, an accepted contribution is linked to the commit and
files where it was included. Corrections add to that history.

A compute contributor needs a well-chosen task. The default work packet names
an open statement, relevant source material, a permitted experiment or proof
attempt, and a stopping condition. Preparing useful packets is mathematical
work; qualified review should precede large-scale distribution. The public
source-packet exporter makes a selected task portable: a contributor chooses
committed files and a question, then uploads the resulting text attachment to
a model outside the repository harness. The packet preserves the source pin
and hashes so that the returned argument can be checked against the same
starting material. Selecting enough context remains the contributor's job.
That review
has not been established for every current task.

Mathematical direction and verification can be contributed separately. Someone
may propose a promising question without running an agent, while another
contributor examines the eventual result. The record attributes both roles.

Returns are reviewed as mathematics or as changes to the research software.  The \emph{mathematics track}
accepts proofs, reductions, computations, counterexamples, corrections,
formalisation, and reproducible failed routes.  The \emph{architecture track}
accepts improvements to agent workflow, navigation, validation,
reproducibility, public experience, governance, and tooling.  An architecture
return never needs a fictitious problem number, and it cannot request
promotion of a mathematical claim.

The return, review decision, and release remain separate records. A useful
correction or design proposal can therefore retain credit even if it does not
lead to a merged theorem.

\section{From an agent claim to mathematical attention}\label{sec:attention}

A returned result needs enough preparation that a specialist can assess it.
The reviewer should be able to recover the statement, reproduce the evidence,
find the relevant prior work, and identify the step that remains uncertain.
For a formal proof, this includes checking that the Lean proposition expresses
the intended mathematics.

The maintainer first checks whether the return is coherent and relevant to the
project. A sufficiently developed result can then be prepared for external
attention. Comparator checks selected formal interfaces; Palomar records
formal results with structured disclosure and an automated editorial filter.
Palomar does not provide human expert review or establish novelty
\cite{palomar}. The Erd\H{o}s Problems community provides a place to discuss
results relevant to its questions, including links to longer arguments
\cite{erdosproblems}.

Independent replay, a resolved objection, and a clear explanation give an
expert specific reasons to spend time on a result. The project should preserve
those contributions alongside the proof, so the next reviewer need not
reconstruct the entire discussion.

\section{Progress, provenance, and credit}\label{sec:credit}
\papersectiontarget{strategy-credit}

Credit records what each contributor did. It can name the author of a theorem,
the person who found a counterexample, the contributor who formalised a proof,
and those who corrected, reviewed, or explained it. Software and computation
receive credit too. The current receipt schema optionally uses the fourteen
CRediT roles, which describe contributions without deciding authorship
\cite{credit}.

Credit travels with the repository. After a maintainer accepts a contribution,
its record links the people and work to the evidence, public files, and commit
where it was included. The public contribution pages are generated from these
committed records, so a reader can follow a credit entry back to the work it
recognises. Later corrections add to the history. The project does not require
ownership of an underlying idea in exchange for including it. Authorship of a
later paper or library contribution is a separate decision based on the work
performed.

This distinction matters when a local lemma is generalised. The original
observation should remain cited, while the mathematician and Lean contributor
who develop the general theorem and library interface receive credit for that
work. A reference back to the motivating problem need not compete with their
authorship.

\subsection*{Later uses of a contribution}

A contribution's importance may become clearer through later work. The record
can link an earlier observation to a theorem that uses it, a correction, or a
new application. These links are claims to inspect and revise, rather than
points on a leaderboard. The original evidence and credit remain available
even when no further use has been found. Financial rewards would require a
separate policy; the present record promises attribution.

\section{Distributed compute without distributed authority}\label{sec:security}
\papersectiontarget{strategy-security}

Volunteer compute has an established precedent.  BOINC packages scientific
jobs for heterogeneous consumer devices, and BOINC Central explicitly aims to
make volunteer computing available to scientists without the resources to
operate their own project \cite{boinc}.  GIMPS shows a mathematical version:
volunteers run a common search program, independent machines confirm a prime,
and discovery credit includes the compute donor, software authors, server
operator, and wider volunteer effort \cite{gimps}.

Agent-assisted research is less uniform than either example.  Most tasks are
not interchangeable work units with a predetermined verifier.  A runner may
change code, propose a conjecture, or misread the problem.  The public system
therefore distributes \emph{search} while keeping authority local to each
evidence type.  Mathematicians and formalisation contributors design or review
the routes; volunteers may execute them with Claude Code, Codex, Cursor,
Antigravity, another open or proprietary harness, or no agent at all.  The
runner is replaceable.  The task packet, evidence boundary, and return record
are the shared protocol.  Returned code is untrusted until reviewed.  Expensive or
privileged continuous-integration jobs must not execute fork code with
repository secrets.  In particular, a GitHub \texttt{pull\_request\_target}
workflow must not check out and execute an untrusted fork; GitHub documents
that combination as a route to secret leakage and repository compromise
\cite{githubsecurity}.  Initial checks should run with read-only permissions,
no secrets, bounded resources, and explicit artefact retention.  Promotion to
trusted infrastructure is a later review decision.

A useful compute return records at least the source commit, model and runner,
tool disclosure, prompt or task packet, wall-clock time, human time, token or
subscription use, and hardware budget where available, together with commands,
changed paths, outputs, validation receipts, failures, and the stopping rule.
These inputs should not be forced into one exchange rate: a mathematician's
hour, an API bill, and a donated graphics processor are different resources.
They can still be disclosed well enough to make a capability comparison
interpretable and to reveal when hidden scaffolding paid most of the cost.

The first public mode may remain deliberately simple: contributors clone the
repository and run their own agents locally.  A later volunteer-compute layer
could distribute signed, immutable task packets and receive result bundles
without granting write access.  It should be built only after task identity,
sandboxing, deduplication, resource budgets, result replay, and abuse handling
have explicit owners.

\section{Related projects}\label{sec:prior}

The contribution process draws on public mathematical collaboration, shared
formal libraries, and volunteer computing. The following comparisons explain
which practices are being reused.

\paragraph{Distributed compute.}
BOINC and GIMPS show that members of the public will donate hardware to a
scientific objective when the client is easy to run, the work is visible, and
credit is legible \cite{boinc,gimps}.  Their tasks are much more mechanically
uniform than open-ended proof research.

\paragraph{Distributed mathematical insight.}
Polymath projects invite participants at different mathematical levels to
share small observations as they occur.  Their rules explicitly welcome
tentative and negative insights, provided they are made clear enough for
others to absorb, and treat the project as collaboration rather than a race
\cite{polymath}.  Polymath supplies a social protocol; it does not supply a
formal proof, computation, and agent-return substrate for every comment.

\paragraph{Distributed formalisation.}
Mathlib uses ordinary fork-and-pull-request practice, human review, and
source-level attribution.  The Carleson formalisation used a public blueprint
and many claimable lemma-sized tasks, with formalisation feeding corrections
back into the informal plan \cite{mathlib,carleson}.  Google DeepMind's
\texttt{formal-conjectures} repository similarly uses one-problem files,
issue assignment, pull requests, metadata, and stable snapshots for a growing
Lean benchmark \cite{formalconjectures}.  These projects demonstrate that
large checked mathematical objects can be made publicly divisible.

\paragraph{Multi-agent formal research.}
Agent Hunt studies bounties, locks, guarded ownership, and collaborative
agents for autoformalisation \cite{agenthunt}.  Lean Atlas uses formal
dependency information to reduce the declarations a person must inspect for
semantic verification \cite{leanatlas}.  LeanMarathon adds a durable proof blueprint, scoped agents, and checked
parallel integration for research-paper formalisation, including prose stored
beside formal statements and checked against their proof dependencies
\cite{leanmarathon}.
These are close precedents for the working environment. The proposal here is
to make its continuing mathematical record and infrastructure jointly
contributable, with evidence and credit attached to accepted changes.

\paragraph{Proof abundance.}
Tao separates problem solving into generation, verification, exposition,
digestion and acceptance, and canonicalisation, and argues that proof
abundance will create bottlenecks between these stages \cite{tao2026}.  The
strategy here treats those bottlenecks as contribution surfaces.  A reviewer
who prevents an overclaim or an expositor who makes a hard step recoverable is
not ancillary to the research pipeline.

The proposed composition joins these precedents.  It combines volunteer
compute, small shared insights, formal task decomposition, agent-native
problem worlds, explicit negative knowledge, untrusted public returns,
role-aware credit, and human review.  The candidate contribution is this
composition and its implementation in one open mathematical corpus.  The
paper reports no controlled comparison showing that the composition increases
discovery rate.

\section{Adding another problem world}\label{sec:add-problem}
\papersectiontarget{strategy-add-problem}

The architecture is intended to hold more than the present eight problems,
but adding a folder is not enough.  A coherent new world needs a canonical
question and primary source; a dated account of its current status; an explicit
public claim boundary; a readable primer or problem note; formal statements
with an informal-faithfulness account where formalisation is appropriate;
known literature, computations, counterexamples, and failed routes; exact open
contribution points; navigation and validation routes; attribution; and a
maintainer or review path.

The public skill for adding a problem separates three states.  A proposal may
begin as a sourced issue.  An incubating formal lane may add checked source
under a problem-owned namespace while stating that it is not yet a reviewed
public claim.  A fully indexed world joins the problem registry, paper corpus,
query routes, return schema, cold-start inventory, and relevant external
crosswalks.  These are distinct transitions because a checked proposition, a
published paper, a reviewed claim, and a Comparator selection are different
facts.

After any of these transitions, consequence propagation inspects the problem
registry, paper corpus, query routes, cold-start inventory, validation roster,
return schema, and external crosswalks. A new problem is complete at its stated
entry level only when every affected consumer has an explicit disposition; a
new directory is not itself a lifecycle receipt.

The last transition is not yet a one-command public operation.  The current
paper-corpus fan-in has a maintainer-owned stage, several checks enumerate the
present roster, and the bounded problem index is already close to its size
limit.  The immediate infrastructure work is therefore to derive rosters from
one public authority, admit an explicit incubating state, split verbose detail
out of the bounded index, and make an absent external crosswalk record a normal
state rather than an error.  Until then, the add-problem skill reports the
couplings instead of concealing them behind a scaffold command.

Arbitrary depth is a design target, not a measured scaling result.  The useful
test is whether a cold agent can retrieve the relevant neighbourhood without
reading every paper or rediscovering an existing failure as the number and
depth of worlds grow.  That test should be repeated whenever the roster or
navigation machinery changes.

\section{Participation and growth}\label{sec:growth}

A public contribution process must work for readers with different interests.
A mathematician needs the exact question and prior results. A compute donor
needs a runnable task and a stopping rule. A reviewer needs the claim and its
evidence. An agent builder needs a fixed starting revision against which to
compare runs. The README and entry tools provide routes for each.

The next useful test is a complete outside contribution: a newcomer chooses a
task, returns reproducible work, receives substantive review, and sees the
accepted result and credit in the repository. That would test the proposed
participation model more directly than clone counts or publicity.

\section{Evaluation}\label{sec:evaluation}

The strategy should be evaluated at each conversion boundary.  Useful initial
measures include:

\begin{itemize}
\item time from a cold clone to a correctly scoped first task;
\item fraction of returned bundles that can be replayed at the stated commit;
\item fraction requiring statement or evidence-class correction;
\item time from return to first substantive review and to adoption decision;
\item repeated-dead-end rate before and after a no-go enters the corpus;
\item reviewer time per accepted result family;
\item contribution diversity across mathematics, computation, software,
      validation, exposition, and review;
\item correction lineage completeness and attribution retention; and
\item change in useful output at controlled compute when navigation or another
      infrastructure component changes.
\end{itemize}

The numerator must remain claim-bounded.  One accepted counterexample may be
more useful than hundreds of generated lemmas.  Generated certificate shards
must not be counted as independent discoveries.  A theorem accepted by Lean
but rejected on intended meaning is not a successful public transition.
Likewise, a strong negative result can improve the corpus even though the
endpoint problem remains open.

Longitudinal model comparisons should use immutable problem snapshots and
held-out variants where possible.  Public hard problems are susceptible to
training-data contamination, and the repository itself will become more
informative over time.  A later model therefore receives both a stronger
model and a richer corpus.  The design should measure these factors
separately rather than presenting every later success as model improvement.

The present work supplies no such evaluation.  It establishes an implemented
case, an explicit protocol, and falsifiable measures for a later study.

\section{What to retain from a research attempt}\label{sec:learning}

Retaining a run means selecting the material another researcher could use.
A transcript alone is rarely enough. The record should explain the question,
the evidence found, and the reason for continuing or abandoning the approach.
Search results and computations need source references; a claimed obstruction
needs its exact hypotheses.

Review can also change the order of a paper. The agent must compare the new
work with the current formal source and literature, including results omitted
from the existing manuscript. The strongest useful result should receive the
clearest statement and the space needed to explain its hard step. Routine
lemmas can remain available in the source without determining the paper's
emphasis.

This review starts from the source tree as well as the paper's selected
results. A contributor may find an omitted theorem, an argument assembled
from existing lemmas, or a useful proof that has not reached the public account.
The work then includes checking that argument, formalising a missing step when
needed, and updating the paper and its reading routes. Such a contribution
should be recorded and credited even when it adds few declarations. Its value
is what another researcher can now understand or use.

A problem with the tools is recorded separately. A query that misses a known
lemma may justify a retrieval change; a recurring proof-check failure may
justify clearer instructions. Before generalising such a lesson, the agent
checks whether it recurs in other relevant cases. The repaired instruction
then applies to later work. It does not alter the status of a mathematical
statement. The public maintenance skill asks the next contributor to replay
the repaired journey: find the relevant instructions, run the command, and
follow its result. A reusable repair includes that check and the lesson in the
owning skill, so it survives beyond the conversation that produced it.

\subsection{Generalising a useful lemma}\label{sec:local-to-general}
\papersectiontarget{strategy-local-to-general}

A lemma developed for one problem may have a more general use. Finding that
use requires further mathematics: identify which hypotheses the proof needs,
look for an existing theorem, and formulate a statement that includes the
original result as a special case. Proposed weakenings need counterexample
checks, and the general proof needs verification. Another application can
show whether the generalisation is worth maintaining.

Figure~\ref{fig:local-to-general} follows the proposed route into a shared
library. A mathematician familiar with the subject and a Lean contributor
would need to judge the statement, proof, and library interface before opening
an upstream discussion.

\begin{figure}[H]
\centering
\begin{tikzpicture}[
  node distance=6mm and 9mm,
  box/.style={paper stage,text width=3.3cm,minimum height=1.65cm},
  gate/.style={box,fill=BurntOrange!7,text width=3.7cm},
  arr/.style={paper arrow}
]
\node[box,fill=MidnightBlue!5] (world)
  {{\scriptsize\sffamily\bfseries LOCAL RESULT}\\[3pt]
   \textbf{Hard-problem world}\\[2pt]\scriptsize proof, no-go, or computation};
\node[gate,right=of world] (digest)
  {{\scriptsize\sffamily\bfseries GATE 1}\\[3pt]
   \textbf{Mathematical digestion}\\[2pt]\scriptsize hypotheses, prior art, falsifiers};
\node[box,fill=TealBlue!5,right=of digest] (generic)
  {{\scriptsize\sffamily\bfseries CANDIDATE}\\[3pt]
   \textbf{Reusable theorem}\\[2pt]\scriptsize natural statement and consumers};
\node[gate,below=of generic] (expert)
  {{\scriptsize\sffamily\bfseries GATE 2}\\[3pt]
   \textbf{Expert stewardship}\\[2pt]\scriptsize meaning, fit, maintenance, credit};
\node[box,fill=ForestGreen!6,left=of expert] (library)
  {{\scriptsize\sffamily\bfseries SHARED OUTPUT}\\[3pt]
   \textbf{Library result}\\[2pt]\scriptsize reused, or declined with reasons};
\draw[arr] (world) -- (digest);
\draw[arr] (digest) -- (generic);
\draw[arr] (generic) -- (expert);
\draw[arr] (expert) -- (library);
\end{tikzpicture}
\caption{The proposed route from a problem-specific lemma to a shared library.
The generalisation and library review are additional work.}
\label{fig:local-to-general}
\end{figure}

Mathlib is the clearest external model for the final stage, and also a warning
against automating it badly.  Its contribution guide welcomes useful
contributions and supplies a public fork-and-pull-request route; it does not
reserve contribution to people with institutional credentials.  It also sets
high standards for generality, integration, maintainability, style, and
documentation, and its current AI policy rejects low-quality unsupervised
agent submissions while requiring AI use to be disclosed \cite{mathlib}.
Accordingly, this project should never direct a compute donor to spray raw
agent lemmas into Mathlib.  A subject and Lean expert must take custody of a
candidate, understand it, decide whether it belongs, bring it to community
standards, and own the discussion.  The public contribution is then theirs in
the roles they performed; the originating problem route remains a provenance
fact, not a competing claim of ownership.

The present prototype does not automate this track and reports no Mathlib
contribution produced by it.  It can preserve local evidence and provenance,
and it can help prepare a candidate.  Canonical status and upstream acceptance
remain external outcomes.  The
\papersectionlink{claim-faithful-publication-systems-paper.pdf}%
{systems-mathloop}{systems paper's mathematical loop} describes the authority
separation underneath this handoff.

The resulting record can connect a local lemma to a later generalisation, an
application, or a correction. It can also retain unsuccessful arguments with
the reason they fail. These records may be useful for retrieval or future
model training. Whether they improve mathematical originality remains an
experimental question.

\section{Possible applications outside mathematics}\label{sec:science}

A later extension could investigate whether the same record-and-review
approach helps simulation or analysis of public scientific data. Physical
experimentation would additionally require domain expertise, calibrated
instruments, uncertainty analysis, controlled facilities, and safety review.
The mathematical prototype supplies no evidence of those capabilities;
autonomous physical experimentation is not a present project goal.

\section{Limits and governance}\label{sec:limits}
\papersectiontarget{strategy-limits}

The prototype had no recorded external user or accepted outside contribution
by 31 August 2026. Setup and review may therefore have problems that an
author-operated rehearsal misses. A clean reproduction or a report of a broken
instruction would be useful evidence about the contribution process.

Decisions also remain concentrated in one maintainer. Outside reviewers and
maintainers would need a part in problem selection, adoption decisions, and
credit disputes. Published correction and appeal procedures would make those
decisions open to challenge.

More agents could increase the review burden without producing worthwhile
mathematics. Evaluation should therefore account for rejected work, reviewer
time, and possible exposure to the public corpus during model training.
Contributors have unequal access to compute and expertise; sharing the
software addresses only part of that difference.

\section{Execution order}\label{sec:execution}

The immediate task is to test the local contribution process with outside
users: finding a question, running the relevant tools, returning evidence, and
receiving review and credit. The README, query tools, return schemas, and
contribution records supply the starting implementation.

External reviewers and maintainers are needed before a substantial increase
in agent throughput. Controlled comparisons can then test changes to retrieval,
models, and compute budgets against fixed source revisions. A hosted volunteer
scheduler, public task market, or cross-domain service would be later work;
none is a reported capability of this release.

\section{Conclusion}\label{sec:conclusion}

The proposal is to share a usable research environment together with its
mathematical results. Each accepted contribution can improve the problem
record or the tools used to investigate it, and its credit remains attached.
A contributor can therefore advance one part of a difficult problem without
reconstructing the entire project.

The public repository implements this contribution path for eight open
Erd\H{o}s problems. Whether it becomes a productive research commons depends on
external use: contributors finding worthwhile tasks, reviewers being able to
check the returns, and later work benefiting from what was retained. Those are
the next outcomes to measure.

\appendix
\section{Public entry routes}\label{app:routes}
\papersectiontarget{strategy-entry-routes}

The public repository is \href{\repobase}{\texttt{wcook04/plectis-lean-erdos249-257}}.
The shortest current routes are:

\begin{center}
\small
\begin{tabular}{@{}p{0.35\linewidth}p{0.54\linewidth}@{}}
\toprule
\papertablehead
Question & Public route \\
\midrule
What is the experiment? & \repolink{README.md}{README.md} \\
Where should an agent begin? & \repolink{AGENTS.override.md}{AGENTS.override.md} \\
How can an agent explain the system? &
\repolink{skills/explain-public-system/SKILL.md}{explain-public-system skill} \\
How can an agent run the coupled research lifecycle? &
\repolink{skills/run-coupled-research-goals/SKILL.md}{coupled-goal skill} \\
How can an agent mine a frontier? &
\repolink{skills/mine-open-problem/SKILL.md}{mine-open-problem skill} \\
How is a bounded Lean change validated? &
\repolink{skills/lean-concurrent-validation/SKILL.md}{lean-concurrent-validation skill} \\
How are downstream consequences reconciled? &
\repolink{skills/propagate-research-consequences/SKILL.md}{propagate-research-consequences skill} \\
How are clone skills installed elsewhere? &
\repolink{skills/install-clone-skills/SKILL.md}{install-clone-skills skill} \\
What is proved and what remains open? & \repolink{docs/RESULTS.md}{docs/RESULTS.md} \\
Which papers and problems exist? & \repolink{docs/papers/README.md}{docs/papers/README.md} \\
How can I contribute mathematics? & \repolink{CONTRIBUTING.md}{CONTRIBUTING.md} and the
research-progress issue form \\
How can I improve the architecture? &
\repolink{docs/research-commons/ARCHITECTURE_CONTRIBUTIONS.md}{architecture contribution guide}
and the architecture-proposal issue form \\
How can an agent prepare a pull request? &
\repolink{skills/submit-pull-request/SKILL.md}{submit-pull-request skill} \\
How can I propose or add another problem? &
\repolink{skills/add-open-problem/SKILL.md}{add-open-problem skill} \\
How is a return validated? &
\repolink{skills/erdos-research-return/SKILL.md}{erdos-research-return skill} and the
\repolink{docs/research-commons/README.md}{research-commons protocol} \\
How is credit recorded? &
\repolink{docs/research-commons/CREDIT_POLICY.md}{credit policy} \\
What do Comparator and Palomar establish? &
\repolink{docs/EXTERNAL_VERIFICATION.md}{Comparator guide} and
\repolink{docs/PALOMAR_QUALIFICATION.md}{Palomar qualification} \\
\bottomrule
\end{tabular}
\end{center}

\section*{Declaration of generative AI use}
Every word of this manuscript was generated by agents based on large language
models operating within Will Cook's private research system for artificial
intelligence. The formal proofs and repository software were likewise drafted
and revised by the agents through that system under Cook's direction. Cook set
the objectives and acceptance criteria and maintains the research infrastructure.
Source-bound mathematical interpretation was performed by the agents, not
independently verified by Cook. Cook assumes responsibility for
the accuracy, interpretation, and presentation of the work. Generative systems
are production tools, not authors, and supply no independent authority.

\begin{thebibliography}{99}

\bibitem{leanmarathon}
Y. Zhang, Y. Sun, T. Suzuki, J. D. Lee, and F. Liu,
\emph{LeanMarathon: Toward Reliable AI Co-Mathematicians through Long-Horizon
Lean Autoformalization}, 2026,
\href{https://arxiv.org/abs/2606.05400}{arXiv:2606.05400}, Sections 2--4.


\bibitem{tao2026}
T. Tao, \emph{Mathematics in the age of AI}, 2026,
\href{https://arxiv.org/abs/2608.16753}{arXiv:2608.16753}.

\bibitem{taolearning2026}
T.~Tao, discussion of learning from studying a mathematical problem,
Mastodon thread, 3 September 2026,
\href{https://mathstodon.xyz/@tao/117208617602946453}{opening post},
accessed 5 September 2026.
\bibitem{taoexploration2026}
T.~Tao, discussion of exploration and insight in the Navier--Stokes
regularity problem, Mastodon thread, 3 September 2026,
\href{https://mathstodon.xyz/@tao/117207849921390904}{opening post};
\href{https://mathstodon.xyz/@tao/117207855800042681}{part 5 on preserving exploration},
accessed 5 September 2026.

\bibitem{boinc}
D. P. Anderson, \emph{BOINC: A Platform for Volunteer Computing}, Journal of
Grid Computing 18 (2020), 99--122,
\href{https://doi.org/10.1007/s10723-019-09497-9}{DOI}.

\bibitem{gimps}
Great Internet Mersenne Prime Search, \emph{GIMPS}, project documentation and
discovery-credit record, \href{https://www.mersenne.org/}{mersenne.org},
accessed August 2026.

\bibitem{polymath}
Polymath Project, \emph{General polymath rules},
\href{https://polymathprojects.org/general-polymath-rules/}{project rules},
accessed August 2026.

\bibitem{mathlib}
Lean community, \emph{Contributing to mathlib},
\href{https://leanprover-community.github.io/contribute/index.html}{contributor guide},
accessed August 2026.

\bibitem{carleson}
L. Becker et al., \emph{A Blueprint for the Formalization of Carleson's
Theorem on Convergence of Fourier Series}, 2025,
\href{https://arxiv.org/abs/2405.06423}{arXiv:2405.06423}.

\bibitem{erdosproblems}
T. F. Bloom, \emph{Erd\H{o}s Problems}, problem database, sources, and forum,
\href{https://www.erdosproblems.com/}{erdosproblems.com},
accessed August 2026.

\bibitem{lean4}
L. de Moura and S. Ullrich, \emph{The Lean 4 Theorem Prover and Programming
Language}, Automated Deduction---CADE 28 (2021), 625--635,
\href{https://doi.org/10.1007/978-3-030-79876-5_37}{DOI}.

\bibitem{formalconjectures}
Google DeepMind, \emph{formal-conjectures},
\href{https://github.com/google-deepmind/formal-conjectures}{software repository},
accessed August 2026.

\bibitem{palomar}
Palomar Registry, \emph{About Palomar} and \emph{Contribution policy},
\href{https://palomar-registry.org/about}{registry documentation} and
\href{https://github.com/PalomarRegistry/PalomarPolicy/blob/main/CONTRIBUTING.md}{submission standard},
accessed August 2026.

\bibitem{agenthunt}
C. E. Brown, C. Kaliszyk, and J. Urban, \emph{Agent Hunt: Bounty Based
Collaborative Autoformalization With LLM Agents}, 2026,
\href{https://arxiv.org/abs/2603.06737}{arXiv:2603.06737}.

\bibitem{leanatlas}
B. Yanahama and A. Sannai, \emph{Lean Atlas: An Integrated Proof Environment
for Scalable Human--AI Collaborative Formalization}, 2026,
\href{https://arxiv.org/abs/2604.16347}{arXiv:2604.16347}.

\bibitem{credit}
NISO, \emph{CRediT: Contributor Roles Taxonomy},
\href{https://credit.niso.org/contributor-roles-defined/}{role definitions},
accessed August 2026.

\bibitem{githubsecurity}
GitHub, \emph{Preventing pwn requests}, GitHub Actions security guidance,
\href{https://docs.github.com/en/actions/reference/security/secure-use}{documentation},
accessed August 2026.

\end{thebibliography}

\end{document}
